\sectionheading{5.7 Java}

Java একটি \textbf{High Level, General Purpose এবং Object-Oriented
Programming Language}। এটি এমনভাবে তৈরি করা হয়েছে যাতে Programmer
তুলনামূলক সহজ ও বোধগম্য Code লিখতে পারে এবং প্রয়োজনীয় Java
Environment থাকলে একই Program বিভিন্ন Platform-এ ব্যবহার করা যায়।

%================================================
% Definition
%================================================
\booksubheading{Java কাকে বলে?}

\begin{definitionbox}

\textbf{\color{darkTeal}সংজ্ঞা}

Java হলো একটি \textbf{High Level, Object-Oriented ও General Purpose
Programming Language}, যা বিভিন্ন ধরনের Software ও Application
তৈরিতে ব্যবহৃত হয়।

\end{definitionbox}

\textbf{সহজভাবে মনে রাখবে—}

\begin{center}

\textbf{
Java
=
High Level
+
Object-Oriented
+
General Purpose
}

\end{center}

%================================================
% History
%================================================
\booksubheading{Java-এর সংক্ষিপ্ত ইতিহাস}

Java-এর ইতিহাসের কয়েকটি তথ্য MCQ-এর জন্য গুরুত্বপূর্ণ—

\begin{compacttwocolpoints}

\item Java-এর প্রধান উদ্ভাবক \textbf{James Gosling}।

\item ভাষাটি \textbf{Sun Microsystems}-এ তৈরি হয়।

\item Java-এর প্রাথমিক নাম ছিল \textbf{Oak}।

\item Java আনুষ্ঠানিকভাবে \textbf{1995 সালে} প্রকাশিত হয়।

\item পরবর্তীতে Sun Microsystems-কে Oracle Corporation অধিগ্রহণ করে।

\end{compacttwocolpoints}

\begin{rememberbox}

\textbf{\color{deepOrange}মনে রাখবে}

\begin{center}

\textbf{
James Gosling
\(\rightarrow\)
Oak
\(\rightarrow\)
Java
\(\rightarrow\)
1995
}

\end{center}

\end{rememberbox}

%================================================
% Features
%================================================
\booksubheading{Java-এর প্রধান বৈশিষ্ট্য}

\begin{tabularx}{\linewidth}{|>{\bfseries}p{0.27\linewidth}|Y|}
\hline

\rowcolor{softBlue}
\multicolumn{1}{|c|}{\textbf{বৈশিষ্ট্য}}
&
\multicolumn{1}{c|}{\textbf{সহজ ব্যাখ্যা}}
\\
\hline

High Level
&
মানুষের জন্য তুলনামূলক সহজে লেখা ও বোঝা যায়।
\\
\hline

Object-Oriented
&
Class ও Object ভিত্তিক Programming সমর্থন করে।
\\
\hline

General Purpose
&
বিভিন্ন ধরনের Software ও Application তৈরিতে ব্যবহার করা যায়।
\\
\hline

Platform Independent
&
প্রয়োজনীয় Java Environment থাকলে বিভিন্ন Platform-এ ব্যবহার করা যায়।
\\
\hline

Secure
&
Program ও Runtime-এর নিরাপত্তার জন্য বিভিন্ন ব্যবস্থা রয়েছে।
\\
\hline

Robust
&
Error Handling ও Memory Management-এর সুবিধা রয়েছে।
\\
\hline

Multithreading
&
একই Program-এ একাধিক কাজের ধারা পরিচালনা করা যায়।
\\
\hline

Portable
&
ভিন্ন System-এ Program স্থানান্তর তুলনামূলক সহজ।
\\
\hline

\end{tabularx}

%================================================
% Platform Independence
%================================================
\booksubheading{Java কেন Platform Independent?}

Java-এর সবচেয়ে গুরুত্বপূর্ণ বৈশিষ্ট্যগুলোর একটি হলো
\textbf{Platform Independence}।

Java Source Code সরাসরি নির্দিষ্ট Processor-এর Machine Code-এ
রূপান্তরিত না হয়ে প্রথমে একটি মধ্যবর্তী Code-এ রূপান্তরিত হয়,
যাকে \textbf{Bytecode} বলা হয়।

\begin{center}
\begin{adjustbox}{max width=0.98\linewidth}
\fcolorbox{darkBlue}{softBlue}{
    \parbox[c][8mm][c]{0.19\linewidth}{
        \centering\textbf{Java\\Source Code}
    }
}
\(\longrightarrow\)
\fcolorbox{mainTeal}{softGreen}{
    \parbox[c][8mm][c]{0.17\linewidth}{
        \centering\textbf{Java\\Compiler}
    }
}
\(\longrightarrow\)
\fcolorbox{darkBlue}{softBlue}{
    \parbox[c][8mm][c]{0.16\linewidth}{
        \centering\textbf{Bytecode}
    }
}
\(\longrightarrow\)
\fcolorbox{mainTeal}{softGreen}{
    \parbox[c][8mm][c]{0.13\linewidth}{
        \centering\textbf{JVM}
    }
}
\(\longrightarrow\)
\fcolorbox{darkBlue}{softBlue}{
    \parbox[c][8mm][c]{0.15\linewidth}{
        \centering\textbf{Execution}
    }
}
\end{adjustbox}
\end{center}

এখানে—

\begin{itemize}

\item \textbf{Source Code:} Programmer-এর লেখা Java Program।

\item \textbf{Java Compiler:} Source Code-কে Bytecode-এ রূপান্তর করে।

\item \textbf{Bytecode:} Platform-neutral Intermediate Code।

\item \textbf{JVM:} Bytecode নির্দিষ্ট System-এ Execute করতে সাহায্য করে।

\end{itemize}

%================================================
% JVM
%================================================
\booksubheading{JVM কী?}

\begin{definitionbox}

\textbf{\color{darkTeal}JVM (Java Virtual Machine)}

Java Bytecode-কে নির্দিষ্ট Computer System-এ Execute করার জন্য
ব্যবহৃত Virtual Machine-কে \textbf{JVM} বলা হয়।

\end{definitionbox}

Windows, Linux বা অন্য Platform-এর জন্য উপযোগী JVM থাকতে পারে,
কিন্তু Java Bytecode একই থাকতে পারে।

এই ধারণার সঙ্গে Java-এর বহুল পরিচিত উক্তি হলো—

\begin{center}

\textbf{Write Once, Run Anywhere (WORA)}

\end{center}

অর্থাৎ প্রয়োজনীয় JVM থাকলে একই Java Bytecode বিভিন্ন Platform-এ
চালানো সম্ভব।

\begin{rememberbox}

\textbf{\color{deepOrange}মনে রাখবে}

Java-এর Platform Independence-এর সঙ্গে
\textbf{Bytecode ও JVM} ঘনিষ্ঠভাবে সম্পর্কিত।

\end{rememberbox}

%================================================
% Compiler and File
%================================================
\booksubheading{Java Compiler, Source File ও Bytecode}

Java Program সাধারণত \textenglish{\ttfamily .java} Extension-যুক্ত
Source File-এ লেখা হয়।

যেমন—

\begin{center}

{\textenglish{\ttfamily Hello.java}}

\end{center}

Java Compiler Source Program Compile করলে Bytecode-যুক্ত
\textenglish{\ttfamily .class} File তৈরি হতে পারে।

\begin{center}

\textbf{
\textenglish{\ttfamily .java}
\(\longrightarrow\)
Compile
\(\longrightarrow\)
\textenglish{\ttfamily .class}
}

\end{center}

\textenglish{\ttfamily .class} File-এর Bytecode JVM-এর সাহায্যে
Execute করা হয়।

%================================================
% JDK JRE JVM
%================================================
\booksubheading{JDK, JRE ও JVM-এর প্রাথমিক ধারণা}

Java-এর ক্ষেত্রে তিনটি গুরুত্বপূর্ণ শব্দ হলো
\textbf{JDK, JRE ও JVM}।

\begin{tabularx}{\linewidth}{|>{\bfseries}p{0.22\linewidth}|Y|}
\hline

\rowcolor{softBlue}
\multicolumn{1}{|c|}{\textbf{নাম}}
&
\multicolumn{1}{c|}{\textbf{কাজ}}
\\
\hline

JDK
&
\textbf{Java Development Kit}। Java Program তৈরি, Compile এবং
Development-এর প্রয়োজনীয় Tool-সমূহের সমষ্টি।
\\
\hline

JRE
&
\textbf{Java Runtime Environment}। Java Program চালানোর জন্য
প্রয়োজনীয় Runtime Environment।
\\
\hline

JVM
&
\textbf{Java Virtual Machine}। Java Bytecode Execute করে।
\\
\hline

\end{tabularx}

\textbf{সহজভাবে—}

\begin{center}

\textbf{
JDK = Development
\qquad
JRE = Runtime
\qquad
JVM = Bytecode Execution
}

\end{center}

%================================================
% Java Example
%================================================
\booksubheading{Java-এর একটি ছোট উদাহরণ}

Screen-এ \textenglish{\ttfamily Hello} প্রদর্শনের একটি সাধারণ
Java Program—

\begin{codebox}
\begin{Verbatim}[fontsize=\small]
class Hello {
    public static void main(String[] args) {
        System.out.println("Hello");
    }
}
\end{Verbatim}
\end{codebox}

\textbf{Output:}

\begin{codebox}
\begin{Verbatim}[fontsize=\small]
Hello
\end{Verbatim}
\end{codebox}

এই পর্যায়ে Program-এর বিস্তারিত Structure মুখস্থ করার প্রয়োজন নেই।
এখানে মূল উদ্দেশ্য হলো Java-তে অর্থবোধক Keyword ও Statement
ব্যবহার করে Program লেখা হয়—এই ধারণাটি বোঝা।

%================================================
% Uses
%================================================
\booksubheading{Java-এর ব্যবহার}

Java বিভিন্ন ধরনের Software Development-এ ব্যবহার করা হয়। যেমন—

\begin{compacttwocolpoints}

\item Desktop Application তৈরিতে।

\item Web ও Server-side Application তৈরিতে।

\item Enterprise Software তৈরিতে।

\item Banking ও Business System-এ।

\item Network-based Application-এ।

\item বিভিন্ন বড় Software System তৈরিতে।

\item Android Application Development-এ ঐতিহাসিকভাবে ব্যাপক ব্যবহার হয়েছে।

\end{compacttwocolpoints}

%================================================
% Advantages
%================================================
\booksubheading{Java-এর সুবিধা}

\begin{compacttwocolpoints}

\item Machine Language-এর তুলনায় Syntax অনেক সহজ।

\item Object-Oriented Programming সমর্থন করে।

\item Platform Independence-এর সুবিধা রয়েছে।

\item বড় ও জটিল Software তৈরিতে উপযোগী।

\item Code পুনর্ব্যবহারের সুযোগ রয়েছে।

\item Security ও Error Handling-এর বিভিন্ন সুবিধা রয়েছে।

\item Standard Library সমৃদ্ধ।

\end{compacttwocolpoints}

%================================================
% Limitations
%================================================
\booksubheading{Java-এর কিছু সীমাবদ্ধতা}

\begin{itemize}

\item Java Program চালানোর জন্য উপযুক্ত Java Environment প্রয়োজন।

\item Low Level Language-এর মতো Hardware-এর ওপর সরাসরি নিয়ন্ত্রণ কম।

\item JVM-এর মাধ্যমে Program চলায় কিছু ক্ষেত্রে Native Machine
Code-এর তুলনায় অতিরিক্ত Resource প্রয়োজন হতে পারে।

\end{itemize}

%================================================
% Java vs Python
%================================================
\booksubheading{Java ও Python-এর প্রাথমিক পার্থক্য}

এখন শুধু ধারণার জন্য—

\begin{tabularx}{\linewidth}{|>{\bfseries}p{0.22\linewidth}|Y|Y|}
\hline

\rowcolor{softBlue}
\multicolumn{1}{|c|}{\textbf{বিষয়}}
&
\multicolumn{1}{c|}{\textbf{Java}}
&
\multicolumn{1}{c|}{\textbf{Python}}
\\
\hline

ভাষার স্তর
&
High Level
&
High Level
\\
\hline

ধরন
&
General Purpose, Object-Oriented
&
General Purpose
\\
\hline

Syntax
&
তুলনামূলক বেশি নিয়মবদ্ধ
&
তুলনামূলক সহজ
\\
\hline

Execution ধারণা
&
Bytecode ও JVM
&
Interpreter-ভিত্তিক ব্যবহার বেশি
\\
\hline

গুরুত্বপূর্ণ ব্যক্তি
&
James Gosling
&
Guido van Rossum
\\
\hline

\end{tabularx}

Python পরবর্তী আলাদা Topic-এ বিস্তারিত আলোচনা করা হবে।

%================================================
% Final Reminder
%================================================
\begin{rememberbox}

\textbf{\color{deepOrange}মনে রাখবে}

\begin{itemize}

\item Java একটি \textbf{High Level Programming Language}।

\item Java-এর প্রধান উদ্ভাবক \textbf{James Gosling}।

\item Java তৈরি হয়েছিল \textbf{Sun Microsystems}-এ।

\item Java-এর প্রাথমিক নাম ছিল \textbf{Oak}।

\item Java আনুষ্ঠানিকভাবে প্রকাশিত হয় \textbf{1995 সালে}।

\item Java একটি \textbf{Object-Oriented ও General Purpose Language}।

\item JVM-এর পূর্ণরূপ \textbf{Java Virtual Machine}।

\item Java Source Code Compile হয়ে \textbf{Bytecode} তৈরি হয়।

\item Java Source File-এর Extension সাধারণত
\textenglish{\ttfamily .java}।

\item Bytecode-যুক্ত File-এর Extension সাধারণত
\textenglish{\ttfamily .class}।

\item \textbf{Write Once, Run Anywhere (WORA)} Java-এর
Platform Independence-এর ধারণা প্রকাশ করে।

\item \textbf{JDK = Development, JRE = Runtime, JVM = Bytecode Execution}।

\end{itemize}

\end{rememberbox}

\small

%================================================
% 5.8 Oracle
%================================================
